\section{Related Work}
\label{sec:related}

\noindent\textbf{Data Profiling.}
The closest line of work to \system\ is \emph{data profiling}, which discovers constraints, such as functional and conditional dependencies and DCs, directly from data instances~\cite{profSurveyFelix2015,Huhtala1998TANE,Fan2009DiscoverCFD,Papenbrock2016HyFD,Bleifuss2017Hydra,chu2013discovering,Pena2022FastDC,icdePena2024}. DCs subsume keys, FDs, and ordering constraints~\cite{chu2013discovering}, which makes them our target formalism (Section~\ref{sec:problem}). Approximate discovery tolerates violations in dirty data~\cite{Pena2019DiscoveryDC,martin2026lima}, user-specific criteria help experts rank discovered rules~\cite{fan2023topk}, and LLMs now propose candidate rules from schemas and samples~\cite{biester2024llmclean,glock2026ledqer,li2026lcats}. Still, these methods reveal what the observed data seems to satisfy, not whether a pattern is required by policy or regulation, so profiling cannot recover constraints that are already violated in the released data, and discovered DCs often include many spurious constraints~\cite{martin2025false}.


\smallskip
\noindent\textbf{Data Cleaning.}
Data cleaning systems use constraints, statistical signals, or both to detect and repair errors~\cite{dataerrorsZiawash,chu2013holistic,holoclean,raha2019,peng2022garf,cleaningDeploy2024}, including rules with ML predicates (REEs)~\cite{bao2024rock} and LLM-based detectors~\cite{ni2025zeroed,zhang2024cleaning}. Constraint-based methods are interpretable and auditable, but they assume that the relevant rules are already available. Data-driven methods reduce the need for manual rule specification, but without an external source of truth, they can struggle to distinguish rare valid values from genuine errors~\cite{cleaningDeploy2024}. Our work obtains such rules from regulatory documents, as DCs over a controlled predicate space, which are executable, auditable, and checkable at scale~\cite{facet2021,liu2024rapidash,youri2024}.

\smallskip
\noindent\textbf{Regulatory NLP.}
A separate line of work uses NLP to analyze regulations and compliance-related documents, for example, by extracting rights and obligations or by checking whether requirement statements comply with legal or policy text~\cite{breaux2006regulatory,kiyavitskaya2008automating,amaral2023nlp,hassani2024rethinking,compliance2025ACMTSEM,masoudifard2024graphrag}. LLMs now extract structured rules from regulations~\cite{guliani2026dejure,hettiarachchi-etal-2025-code,kabzhan2025semantic} and translate financial regulations into compliance code~\cite{li2025compliance2code}, although they still struggle to apply long sets of interdependent rules~\cite{yu2026ruleworld,soni2026tam}. These approaches remain text-centered or target a fixed application model, whereas we connect regulatory statements to a relational schema and translate them into constraints whose violations can be checked in the data.

\smallskip
\noindent\textbf{Text-to-SQL.}
Text-to-SQL is also related, since it maps natural-language input to a SQL query over a database schema~\cite{yu2018spider,picard2021,resdsql2023,pourreza2023din,dail2025,llmTEXT2SQLSurvey2025,redsql2025,XiYanSQL,alrashed2026verification,rabonato2026portuguese}. However, Text-to-SQL answers a user question on the current instance, whereas \system\ produces declarative invariants that must hold for all valid instances; beyond linking text to schema elements, it must also identify scope conditions and violation semantics. Closer in output, LLMs generate OCL constraints from specifications of UML models~\cite{abukhalaf2023ocl,attarwala2026ocl}, and Doc2DB systems verify extracted relations against integrity constraints~\cite{zhang2026doc2db}; in both cases, the target schema is designed together with the constraints, whereas \system\ grounds regulatory text in an independently published dataset.
